% =====================================================================
\section{Introduction}\label{sec:intro}
% =====================================================================

Many systems in the applied sciences are governed by evolution
equations that encapsulate the balance between energy storage and
dissipation. This kind of dissipative dynamics underlies, for
instance, the theory of plasticity~\cite{lubliner:1990,hanreddy2013},
viscoelasticity~\cite{tschoegl1989}, damage and
delamination~\cite{lemaitre2005uh}, diffusion in porous
media~\cite{bear1972}, frictional contact~\cite{wriggers2006},
phase transformation in polycrystals~\cite{balljames1987,mielke2002}
and ferromagnetism~\cite{desimone2002,brokate1996}. In all of these
settings the governing equations account for the energetic mechanisms
that take place during the process and ensure that fluxes balance the
accumulation of energy. Remarkably, thermodynamics provides a common
abstract structure that simplifies their analysis and reveals the
essential ingredients of a consistent theory~\cite{germain1983,
halphen1975generalized}.

A particularly rich subclass consists of \emph{rate-independent}
phenomena, whose behaviour is invariant under time
reparametrisation. We refer to the monographs of Mielke and
Roub{\'i}{\v c}ek~\cite{MielkeRoubicek}, of Visintin~\cite{visintin1994}
and of Krasnosel'ski{\u\i} and Pokrovski{\u\i}~\cite{Krasnoselskii}, and
to the foundational work of Mielke and Theil~\cite{MielkeTheil}, for
comprehensive accounts of the mathematical structure of such problems. As shown there,
rate-independent systems possess a variational structure that leads to
powerful analytical results and robust approximation schemes, valid
even in non-smooth settings. Using rate-independent, dissipative model
equations as a template, it becomes easier to propose new models,
irrespective of the field of application, knowing beforehand that the
resulting governing equations respect the most fundamental balance
principles.

Rate-independent models have proven extremely successful in mechanics
and, in particular, in elastoplasticity, fracture and damage
theories~\cite{francfort1998,dalmaso2006,carstensen2002uh}. However,
the framework is far more general, and its transfer to other
disciplines may open fruitful avenues for modelling and analysis.
The present work is motivated by one such transfer: a companion
paper develops a rate-independent model of \emph{epigenetic}
evolution, in which a scalar internal variable encodes the state of
chromatin between a healthy and a damaged phenotype, and an external
oxygen signal plays the role of the loading. Epigenetic marks such as
DNA methylation are heritable, long-lived, and yet actively
reversible~\cite{Bird2002}, and their dependence on the cellular
microenvironment---for example, the impairment of oxygen-dependent
demethylation enzymes under hypoxia~\cite{Thienpont2016}---produces
memory and hysteresis phenomena that are naturally captured by
rate-independent evolution. The biological development is deferred to
the companion paper; here we concentrate on the mathematical framework
that underpins it, which is of independent interest and applicable well
beyond the motivating example.

\paragraph{Scope and contributions.}
Following the abstract set-up of~\cite{MielkeRoubicek}, we specify a
rate-independent system by a triple $(\Q, E, \Psi)$ consisting of a
state space $\Q\subseteq\R^n$, a stored-energy functional $E$
depending on the state and on a prescribed loading path $S(t)$, and a
convex, $1$-homogeneous dissipation potential $\Psi$. The
contributions of this paper are the following.
\begin{enumerate}[leftmargin=2em,itemsep=2pt]
  \item We show that, once the triple is fixed and an energetic
  principle is postulated, the governing equations are determined
  systematically: we derive both the \emph{energetic formulation}
  (global stability and energy balance,
  \cref{def:energetic}) and the equivalent \emph{subdifferential} or
  Biot inclusion (\cref{sec:subdiff}), and we make precise the
  role of the $1$-homogeneity of $\Psi$ through a saturation identity
  (\cref{lem:saturation}).
  \item We give a self-contained thermodynamic reading of the
  framework (\cref{sec:thermo}): the energy balance is the first law,
  the non-negativity of the dissipation rate is the second law, and
  the saturation identity expresses a principle of \emph{minimal}
  (economical) dissipation. Over closed loading cycles this yields a
  purely geometric characterisation of the hysteretic dissipation.
  \item We collect and, where necessary, complete the well-posedness
  theory (\cref{sec:math}): properties of the induced dissipation
  distance, existence of energetic solutions, uniqueness under uniform
  convexity, and existence and (scalar) uniqueness of
  balanced-viscosity solutions obtained by vanishing viscosity, which
  resolve the non-uniqueness of the energetic formulation at jumps.
  \item We analyse a variational time integrator based on incremental
  minimisation (\cref{sec:numerical}): we prove discrete stability and
  a discrete energy inequality, convergence to energetic solutions,
  local consistency of orders one and two under weak and strong
  regularity, and a global energy-consistency estimate. In the scalar
  case the incremental problem reduces to a two-sided return map.
  \item We validate the theory on an abstract scalar double-well
  benchmark driven by a bounded, saturating loading
  (\cref{sec:experiments}), reproducing the analytically known
  quasi-static branches, the reversible and catastrophic hysteresis
  loops, and the predicted first-order energy consistency.
\end{enumerate}

\paragraph{Outline.}
\Cref{sec:framework} introduces the framework and the governing
equations. \Cref{sec:thermo} develops the thermodynamic
interpretation. \Cref{sec:math} presents the well-posedness results,
and \cref{sec:numerical} the variational numerical method and its
analysis. \Cref{sec:experiments} reports the numerical experiments,
and \cref{sec:conclusions} draws conclusions and outlines future work.
To keep the exposition readable, all proofs are collected in the
appendices.
